% from ega3-1-fr.tex
\oldpage[III]{84}
C'est un cas particulier d'une suite exacte de termes de bas degré de la
suite spectrale de Künneth (M, XVII, 5.2 a) et G, I, 5.5.2)~; on peut le
démontrer directement de la façon suivante. Considérons $K_0$ et $K_1$ comme
des complexes de chaînes (nuls en dimensions $\neq0$ et $\neq1$
respectivement)~; on a alors une suite exacte de complexes
\[
0\longrightarrow K_0\otimes L_\bullet
\longrightarrow K_\bullet\otimes L_\bullet
\longrightarrow K_1\otimes L_\bullet\longrightarrow0
\]
à laquelle nous pouvons appliquer la suite exacte d'homologie
\[
\cdots\longrightarrow H_{p+1}(K_1\otimes L_\bullet)
\xrightarrow{\partial}H_p(K_0\otimes L_\bullet)
\longrightarrow H_p(K_\bullet\otimes L_\bullet)
\longrightarrow H_p(K_1\otimes L_\bullet)
\xrightarrow{\partial}H_{p-1}(K_0\otimes L_\bullet)
\longrightarrow\cdots.
\]
Mais il est évident que
$H_p(K_0\otimes L_\bullet)=K_0\otimes H_p(L_\bullet)$ et
$H_p(K_1\otimes L_\bullet)=K_1\otimes H_{p-1}(L_\bullet)$ pour tout $p$~; en
outre, on vérifie aussitôt que l'opérateur
$\partial:K_1\otimes H_p(L_\bullet)\to K_0\otimes H_p(L_\bullet)$ n'est autre
que $d_1\otimes1$~; le lemme résulte donc de la suite exacte précédente et de
la définition de $H_0(K_\bullet\otimes H_p(L_\bullet))$ et de
$H_1(K_\bullet\otimes H_{p-1}(L_\bullet))$.

Ce lemme étant établi, la fin de la démonstration de
(\hyperref[III.1.1.4-fr]{1.1.4}) est immédiate~: l'hypothèse de récurrence et
ce lemme (\hyperref[III.1.1.4.1-fr]{1.1.4.1}) donnent
$H_p(K_\bullet\otimes L_\bullet)=0$ pour $p\geq2$~; en outre, si on prouve que
$H_1(K_\bullet,H_0(L_\bullet))=0$, on déduira aussi
$H_1(K_\bullet\otimes L_\bullet)=0$ du lemme
(\hyperref[III.1.1.4.1-fr]{1.1.4.1})~; mais par définition,
$H_1(K_\bullet,H_0(L_\bullet))$ n'est autre que le noyau de l'homothétie
$z\mapsto f_rz$ dans $M_{r-1}$, et comme par hypothèse ce noyau est nul, cela
achève la démonstration.
\end{proof}

\begin{env}[1.1.5]
\label{III.1.1.5-fr}
Soit $\mathbf g=(g_i)_{1\leq i\leq r}$ une seconde suite de $r$ éléments de
$A$, et posons $\mathbf f\mathbf g=(f_ig_i)_{1\leq i\leq r}$. On peut définir
un homomorphisme canonique de complexes
\begin{equation}
\label{III.1.1.5.1-fr}
\tag{1.1.5.1}
\varphi_{\mathbf g}:K_\bullet(\mathbf f\mathbf g)\longrightarrow
K_\bullet(\mathbf f)
\end{equation}
comme l'extension canonique à l'algèbre extérieure $\bigwedge(A^r)$ de
l'application $A$-linéaire
$(x_1,\ldots,x_r)\mapsto(g_1x_1,\ldots,g_rx_r)$ de $A^r$ dans lui-même. Pour
voir qu'on a bien un homomorphisme de complexes, il suffit de remarquer, de
façon générale, que si $u:E\to F$ est une application $A$-linéaire,
$\mathbf x\in F^\vee$ et $\mathbf y={}^tu(\mathbf x)\in E^\vee$, alors on a
la formule
\begin{equation}
\label{III.1.1.5.2-fr}
\tag{1.1.5.2}
(\bigwedge u)\circ i_{\mathbf y}=i_{\mathbf x}\circ(\bigwedge u)~;
\end{equation}
en effet, les deux membres sont des antidérivations de $\bigwedge F$, et il
suffit de vérifier qu'ils coïncident dans $F$, ce qui résulte aussitôt des
définitions.

Lorsqu'on identifie $K_\bullet(\mathbf f)$ au produit tensoriel des
$K_\bullet(f_i)$ (\hyperref[III.1.1.1-fr]{1.1.1}), $\varphi_{\mathbf g}$ est
le produit tensoriel des $\varphi_{g_i}$, où $\varphi_{g_i}$ se réduit à
l'identité en degré $0$ et à la multiplication par $g_i$ en degré $1$.
\end{env}

\begin{env}[1.1.6]
\label{III.1.1.6-fr}
En particulier, pour tout couple d'entiers $m,n$ tels que $0\leq n\leq m$,
on a des homomorphismes de complexes
\begin{equation}
\label{III.1.1.6.1-fr}
\tag{1.1.6.1}
\varphi_{\mathbf f^{m-n}}:K_\bullet(\mathbf f^m)\longrightarrow
K_\bullet(\mathbf f^n)
\end{equation}
et par suite des homomorphismes
\begin{align}
\label{III.1.1.6.2-fr}
\tag{1.1.6.2}
\varphi_{\mathbf f^{m-n}} &:K^\bullet(\mathbf f^n,M)\longrightarrow
K^\bullet(\mathbf f^m,M),\\
\label{III.1.1.6.3-fr}
\tag{1.1.6.3}
\varphi_{\mathbf f^{m-n}} &:H^\bullet(\mathbf f^n,M)\longrightarrow
H^\bullet(\mathbf f^m,M).
\end{align}

